A precondition for either margin of leakage is that carbon costs reach buyers as a relative-price increase on regulated suppliers. We document this precondition on the Belgian panel under three identification strategies and use the resulting magnitudes to anchor the CES mapping in \S\ref{sec:domestic_within} and to interpret the across-country results in \S\ref{sec:international}.

\paragraph{Strategy 1: monthly panel local projections on the Känzig structural carbon-policy shock.} On a 2005m1--2019m12 monthly NACE 4-digit PPI panel, regressing $\log \text{PPI}_{s,m+h} - \log \text{PPI}_{s,m-1}$ on the interaction of the \citet{kanzig2023unequal} carbon-policy structural shock $\text{CPShock}_m$ with the sector's pre-shock cost-share intensity $\text{intensity}_s^{\text{base}}$, with sector and year-month fixed effects, gives the impulse response shown in Figure~\ref{fig:ppi_lp_kanzig_style}. The IRF builds gradually from $\widehat{\gamma}_0 \approx 0.17$ on impact, crosses unity around $h = 12$ months ($\widehat{\gamma}_{12} = 1.22$, s.e.\ $0.22$), and peaks at $\widehat{\gamma}_{18} \approx 1.76$ before fading slowly through $h = 36$. Standard errors are clustered on NACE 4-digit. The IRF shape (modest on impact, peak around $h = 15$--$18$, slow fade) is the Belgian-PPI analog of the aggregate-HICP IRF in \citet{kanzig2023unequal}.

\begin{figure}[H]
    \centering
    \includegraphics[width=0.7\textwidth]{phase3_ppi_lp_kanzig_style.pdf}
    \caption{Impulse response of NACE 4-digit log PPI to a unit Känzig structural carbon-policy shock interacted with pre-shock sector cost-share intensity. Sample: 2005m1--2019m12 Belgian NACE 4-digit monthly panel, all sectors. Sector and year-month FE; cluster-robust standard errors at NACE 4-digit. Solid line is the point estimate; dark shaded band is 68\% CI; light shaded band is 90\% CI. Horizon in months on the x-axis.}
    \label{fig:ppi_lp_kanzig_style}
\end{figure}

\paragraph{Strategy 2: regulated-vs-unregulated event study at NACE 4-digit, with sector-specific linear trends.} Following \citet{colmeretal2023}, we estimate
\begin{equation}
    \log \text{PPI}_{s,t} \;=\; \sum_\tau \beta_\tau\,\mathbf{1}[s \in \text{regulated}]\,\mathbf{1}[\text{year} = \tau] \;+\; \alpha_s \;+\; \delta_t \;+\; \theta_s \cdot t \;+\; \varepsilon_{s,t}
    \label{eq:passthrough_strategy2}
\end{equation}
on the 2001--2022 annual NACE 4-digit panel with sector-specific linear trends $\theta_s \cdot t$. Figure~\ref{fig:cmdj_ppi_eventstudy} reports the year-by-year $\widehat{\beta}_\tau$ on the Belgian panel, mimicking \citet{CosterMejeanDiGiovanni2024} Figure 1 on French INSEE PPI data. Phase-aggregated coefficients (CMdG-tight treatment of seven NACE 2-digit codes, the same set as in \citet{CosterMejeanDiGiovanni2024}'s Table A.5):
\[
    \widehat{\beta}_{\text{Phase 2}} \;=\; +0.094\,(***), \qquad \widehat{\beta}_{\text{Phase 4}} \;=\; +0.113\,(**).
\]
Phase 1 and Phase 3 coefficients collapse to insignificance once sector trends are absorbed; we report Phase 2 and Phase 4 because those are the periods where regulated PPI rises differentially after netting out pre-existing sector-level dynamics.

\begin{figure}[H]
    \centering
    \includegraphics[width=0.75\textwidth]{phase1_figure1_cmdj_style.png}
    \caption{Belgian replication of \citet{CosterMejeanDiGiovanni2024} Figure 1: year-by-year coefficients $\widehat{\beta}_\tau$ on $\mathbf{1}[s \in \text{regulated}] \times \mathbf{1}[\text{year} = \tau]$ from equation~(\ref{eq:passthrough_strategy2}), 2001--2022 Belgian NACE 4-digit annual PPI panel with sector-specific linear trends. Vertical bars are 95\% pointwise confidence intervals. Shaded bands denote ETS Phases 1, 2, and 3.}
    \label{fig:cmdj_ppi_eventstudy}
\end{figure}

\paragraph{Strategy 3: descriptive OLS with sector-specific trends.} The annual exposure-share elasticity is small ($+0.07$ in S6 with sector trends) and not significantly different from zero on a single-shock specification --- consistent with the literature \citep[e.g.,][]{martin2014} that finds firm-level pass-through dispersed and noisy in the cross-section. The strongest piece of identifiable evidence remains Strategy 1.

\paragraph{Anchor for the rest of the paper.} For the elasticity bounds in \S\ref{sec:domestic_within} we use the CMdG-tight estimate as the average post-2015 regulated-vs-unregulated log-price differential, magnitude $\approx 0.10$ (\ie about 10 log-points cumulative over Phase~2 and Phase~4). At the right tail, sector-level pass-through is consistent with full pass-through under Strategy 1; we use $\rho = 1$ as an upper bound and $\rho = 0.5$ as a lower bound on the buyer-side relative-price signal.
